% Encoder-ablation variant descriptions (no em dashes).
\paragraph{Encoder variants.}
All three variants share the same tree-derived spike code: continuous features are
converted to spikes at the decision thresholds (split points) harvested from the
trained gradient-boosted tree, so a channel fires when a feature value crosses one of
the boundaries the tree found discriminative. The variants differ only in how many
channels (the population budget) each feature receives, which isolates the contribution
of importance-guided allocation from the choice of thresholds.
\begin{itemize}
  \item \textbf{Tree-guided (\texttt{tree}).} The production encoder. Each feature that
  owns decision edges is given a channel budget proportional to its GBDT importance, so
  informative features are encoded at high resolution (many thresholds) while weak
  features receive a single channel or are dropped. The spike budget is concentrated
  where it carries class information.
  \item \textbf{Uniform (\texttt{uniform}).} Every feature that owns decision edges
  receives exactly one channel, regardless of importance. This keeps the same feature
  set and threshold code as \texttt{tree} but removes the importance weighting, so it
  isolates the effect of the allocation policy by itself.
  \item \textbf{All-channel (\texttt{all}).} Every feature that owns decision edges
  receives the maximum channel budget $M_{\max}$. This is the dense upper bound, an
  encoder that spends spikes on every threshold of every feature with no pruning.
\end{itemize}
Comparing the three separates \emph{which} thresholds are used (shared) from \emph{how
much} spike budget each feature is given (varied). The all-channel encoder roughly
doubles the SOP count relative to the tree-guided encoder for at most $1.7$ points of
macro-F1, and the uniform encoder matches tree-guided accuracy at a similar cost,
confirming that the importance-guided channel budget, not the raw threshold set, is the
mechanism behind the efficiency of Tree-to-Spike.
